\centerline{\bf \Huge Признаки делимости}

\begin{flushright}
        \scriptsize{Именно математика даёт надёжнейшие правила: тому, кто им следует, не опасен обман чувств\\ Леонард Эйлер}
\end{flushright}

\problem{\textbf{1}} Запишем подряд цифры от 1 до 9, получим число 123456789. Простое оно или составное? Изменится ли ответ в задаче, если каким-то образом поменять порядок цифр в этом числе?

\problem{\textbf{2}} Запишите несколько раз подряд число 2013 так, чтобы получившееся число делилось на 9. Ответ объясните.

\problem{\textbf{3}} Найдите числа, которые делятся а) на 9; б) на 4 и заключены между 123456 и 123495.

\problem{\textbf{4}} Напишите наименьшее одиннадцатизначное число, которое делится на 11.

\problem{\textbf{5}} Придумайте а) семизначное; б) десятизначное число, которое делится на 2, 3 и 11, но не
делится на 5 и 9.\\

\problem{\textbf{6}} К числу 15 припишите слева и справа по одной цифре так, чтобы полученное число делилось на 15. Перечислите все возможные варианты.

\problem{\textbf{7}} Чтобы открыть сейф, нужно ввести код — семизначное число, состоящее из двоек и троек.
Сейф откроется, если двоек в коде больше, чем троек, а сам код делится и на 3, и на 4. Какой
код может открыть сейф?

\problem{\textbf{8}} У Джерри есть девять карточек с цифрами от 1 до 9. Он выкладывает их в ряд, образуя девятизначное число. Том выписывает на доску все 8 двузначных
чисел, образованных соседними цифрами (например, для числа 789456123 это числа 78, 89, 94,
45, 56, 61, 12, 23). За каждое двузначное число, делящееся на 9, Том отдает Джерри кусочек
сыра. Какое наибольшее количество кусочков сыра может получить Джерри?\\

\problem{\textbf{9}} Замените звездочки в записи числа 72*4* цифрами так, чтобы это число делилось на 45.
Сколько решений имеет эта задача? Выпишите все решения.

\problem{\textbf{10}} Известно, что число 55*44** делится на 72. Какие цифры могут стоять вместо звёздочек? Перечислите все возможные варианты.

